\documentclass{article}
\usepackage[UTF8]{ctex}
\usepackage{amsmath,amssymb,amsfonts}
\usepackage{enumitem}
\usepackage{geometry}
\geometry{a4paper, margin=1in}
% === 错题本排版优化（不参与题目标识） ===
\setlength{\parindent}{0pt}
\setlength{\parskip}{0.35em plus 0.10em minus 0.08em}
\setlength{\emergencystretch}{3em}
\setlist{itemsep=0.25em, topsep=0.40em, parsep=0pt}
\allowdisplaybreaks[2]
% === 错题本排版优化结束 ===
\begin{document}
\section*{中国石油大学(华东)2025年高等代数真题}
\section*{26．中国石油大学（华东）2025年高等代数试卷（回忆版）}

\subsection{题目1}
1．已知 \(\lambda^{4}-3 \lambda^{2}+a \lambda-b-2=0\) 有至少 2 个有理根，且其中一个大于 1 ，求 \(a, b\) ．

\subsection{题目2}
2．求解关于 \(x\) 的方程

\[
\left|\begin{array}{cccccc}
x & a_{2} & a_{3} & \cdots & a_{n-1} & 1 \\
a_{1} & x & a_{3} & \cdots & a_{n-1} & 1 \\
a_{1} & a_{2} & x & \cdots & a_{n-1} & 1 \\
\vdots & \vdots & \vdots & & \vdots & \vdots \\
a_{1} & a_{2} & a_{3} & \cdots & x & 1 \\
a_{1} & a_{2} & a_{3} & \cdots & a_{n-1} & 1
\end{array}\right|=0
\]

的根。

\subsection{题目3}
3．讨论线性方程组

\[
\left\{\begin{array}{l}
\lambda x_{1}+x_{2}+x_{3}+x_{4}=1 \\
x_{1}+\lambda x_{2}+x_{3}+x_{4}=\lambda \\
x_{1}+x_{2}+\lambda x_{3}+x_{4}=\lambda^{2} \\
x_{1}+x_{2}+x_{3}+\lambda x_{4}=\lambda^{3}
\end{array}\right.
\]

当 \(\lambda\) 取何值时有唯一解？无穷多解？无解？并在有解时求出解。

\subsection{题目4}
4．\(V\) 是 \(n\) 维线性空间，\(V_{1}\) 是其子空间，且 \(\operatorname{dim} V_{1} \geq \frac{n}{2}\) 。证明：

\begin{itemize}
  \item[（1）] 存在 \(W_{1}, W_{2}\) 是 \(V\) 的子空间，满足 \(V=V_{1} \oplus W_{1}=V_{1} \oplus W_{2}\) ，且 \(W_{1} \cap W_{2}=\{0\}\) 。
  \item[（2）] 当 \(\operatorname{dim} V_{1}<\frac{n}{2}\) 时，上述结论是否成立？并说明理由。
\end{itemize}

\subsection{题目5}
5．证明：求导变换 \(\sigma\) 的非零不变子空间中必有 1 ，进一步证明：\(\sigma\) 的 \(k\) 维不变子空间一定是由\(1, x, \cdots, x^{k-1}\) 生成的子空间．

\subsection{题目6}
6．\(A\) 是 3 阶正交矩阵，且 \(|A|=1\) 。

\begin{itemize}
  \item[（1）] 证明： 1 是 \(A\) 的特征值．
  \item[（2）] 证明：\(A\) 的特征多项式可表示为 \(f(\lambda)=\lambda^{3}-a \lambda^{2}+a \lambda-1, a\) 为一个实数．
  \item[（3）] 若 \(A\) 的特征值全为实数，且 \(|A+E| \neq 0\) ，证明：\(A^{\prime}=A^{2}-3 A+3 E\) 。
\end{itemize}

\subsection{题目7}
7．二次型 \(f\left(x_{1}, x_{2}, x_{3}\right)=5 x_{1}^{2}+5 x_{2}^{2}+c x_{3}^{2}+2 x_{1} x_{2}+6 x_{1} x_{3}-6 x_{2} x_{3}\) 秩为 2 。

\begin{itemize}
  \item[（1）] 求 \(c\) 。
  \item[（2）] 求正交线性替换将其化为标准型。
  \item[（3）] \(f=1\) 时代表什么二次曲面？
\end{itemize}

\subsection{题目8}
8．解答如下问题：

\begin{itemize}
  \item[（1）] 证明：\(r(A B)+r(B C) \leq r(A B C)+r(B)\) ．
  \item[（2）] \(A=\left(\begin{array}{cccc}a_{1}^{2}-1 & a_{1} a_{2} & \cdots & a_{1} a_{n} \\ a_{2} a_{1} & a_{2}^{2}-1 & \cdots & a_{2} a_{n} \\ \vdots & & & \vdots \\ a_{n} a_{1} & a_{n} a_{2} & \cdots & a_{n}^{2}-1\end{array}\right)\) ，求证：\(r(A) \geq n-1\) ，并说明取等号的充要条件．
\end{itemize}

\subsection{题目9}
9．\(A=\left(\begin{array}{llll}2 & 4 & 0 & 0 \\ 1 & 2 & 0 & 0 \\ 0 & 0 & 2 & 0 \\ 0 & 0 & 4 & 2\end{array}\right)\) ，求 \(A^{n}\) ．
\end{document}
